\documentclass[nonacm]{acmart}

\AtBeginDocument{%
  \providecommand\BibTeX{{%
    Bib\TeX}}}

\setcopyright{none}

\usepackage{array}
\usepackage{mathtools}
\usepackage{stmaryrd}
\usepackage{mathpartir}
\usepackage{listings}
\usepackage{xcolor}

\lstset{
  basicstyle=\ttfamily\small,
  breaklines=true,
  columns=fullflexible,
  keepspaces=true,
  frame=single,
  showstringspaces=false
}

\newcommand{\ctx}{\mathrel{\mathsf{ctx}}}
\newcommand{\FV}{\mathsf{FV}}
\newcommand{\dom}{\mathsf{dom}}

\begin{document}

\title{Local Type Definitions in System F$_\omega$}

\author{Runbang Yan}
\email{rbyan25@stu.pku.edu.cn}
\affiliation{%
  \institution{School of Computer Science, Peking University}
  \city{Beijing}
  \country{China}
}

\renewcommand{\shortauthors}{Yan}

\maketitle

\section{Extension}

We extend System F$_\omega$ with transparent local type definitions:
\[
t ::= \cdots \mid \mathbf{let}\ X = T\ \mathbf{in}\ t
\qquad
\Gamma ::= \cdots \mid \Gamma, X :: K = T .
\]
The context entry $X :: K = T$ binds $X$ as a local abbreviation for $T$.
It is not an abstract type variable: inside the scope of the binding, $X$ is definitionally equal to $T$.
All substitutions below are capture-avoiding.

\section{Context Formation}

The new context formation rule is:
\[
\inferrule*[Right=\textsc{C-TDef}]
  {
    \Gamma \ctx
    \\
    \Gamma \vdash T :: K
    \\
    X \notin \dom(\Gamma)
  }
  {
    \Gamma, X :: K = T \ctx
  }.
\]
The premise $\Gamma \vdash T :: K$ prevents recursive type abbreviations such as
$X :: K = X$.
Thus local type definitions do not add unrestricted type-level recursion.

\section{Kinding}

A defined type variable can be used with its declared kind:
\[
\inferrule*[Right=\textsc{K-TDefVar}]
  {
    X :: K = T \in \Gamma
  }
  {
    \Gamma \vdash X :: K
  }.
\]
All existing kinding rules are unchanged.
For example, if $\Gamma, X :: * = \mathsf{Nat}$ is well formed, then
\[
\Gamma, X :: * = \mathsf{Nat} \vdash X :: *.
\]

\section{Type Equivalence}

The key rule is that a local type definition may be unfolded:
\[
\inferrule*[Right=\textsc{Q-TDef}]
  {
    X :: K = T \in \Gamma
  }
  {
    \Gamma \vdash X \equiv T :: K
  }.
\]
Together with symmetry and transitivity of type equivalence, this also gives
$\Gamma \vdash T \equiv X :: K$.

The ordinary congruence rules of System F$_\omega$ remain valid.  In particular:
\[
\inferrule*[Right=\textsc{Q-Arrow}]
  {
    \Gamma \vdash S_1 \equiv T_1 :: *
    \\
    \Gamma \vdash S_2 \equiv T_2 :: *
  }
  {
    \Gamma \vdash S_1 \to S_2 \equiv T_1 \to T_2 :: *
  }
\]
\[
\inferrule*[Right=\textsc{Q-App}]
  {
    \Gamma \vdash S_1 \equiv T_1 :: K_{11} \Rightarrow K_{12}
    \\
    \Gamma \vdash S_2 \equiv T_2 :: K_{11}
  }
  {
    \Gamma \vdash S_1\ S_2 \equiv T_1\ T_2 :: K_{12}
  }
\]
\[
\inferrule*[Right=\textsc{Q-Abs}]
  {
    \Gamma, X :: K_1 \vdash S_2 \equiv T_2 :: K_2
  }
  {
    \Gamma \vdash
    \lambda X :: K_1.\ S_2
    \equiv
    \lambda X :: K_1.\ T_2
    :: K_1 \Rightarrow K_2
  }.
\]
The beta rule for type-level application is unchanged:
\[
\inferrule*[Right=\textsc{Q-Beta}]
  {
    \Gamma, X :: K_{11} \vdash T_{12} :: K_{12}
    \\
    \Gamma \vdash T_2 :: K_{11}
  }
  {
    \Gamma \vdash
    (\lambda X :: K_{11}.\ T_{12})\ T_2
    \equiv
    [X \mapsto T_2]T_{12}
    :: K_{12}
  }.
\]

\section{Typing}

The conversion rule can now use the extended equivalence relation:
\[
\inferrule*[Right=\textsc{T-Eq}]
  {
    \Gamma \vdash t : S
    \\
    \Gamma \vdash S \equiv T :: *
    \\
    \Gamma \vdash T :: *
  }
  {
    \Gamma \vdash t : T
  }.
\]

The typing rule for a local type definition is:
\[
\inferrule*[Right=\textsc{T-TLet}]
  {
    \Gamma \vdash T :: K
    \\
    \Gamma, X :: K = T \vdash t : S
  }
  {
    \Gamma \vdash
    \mathbf{let}\ X = T\ \mathbf{in}\ t
    :
    [X \mapsto T]S
  }.
\]
The result type substitutes away the local name $X$, so $X$ cannot escape its scope.
The usual well-formedness and type-substitution lemmas justify that
$\Gamma \vdash [X \mapsto T]S :: *$.

\section{Evaluation}

Since local type definitions are transparent compile-time abbreviations, evaluation simply substitutes the definition into the body:
\[
\inferrule*[Right=\textsc{E-TLet}]
  { }
  {
    \mathbf{let}\ X = T\ \mathbf{in}\ t
    \longrightarrow
    [X \mapsto T]t
  }.
\]
The substitution acts on type annotations, type applications, package annotations, and all other type occurrences inside the term.

For the example from the slides:
\[
\begin{aligned}
&\mathbf{let}\ X = \mathsf{Nat}\ \mathbf{in}\ 
  (\lambda x:X.\ \mathsf{succ}\ x)\ \mathsf{zero}
\\
&\qquad\longrightarrow
  (\lambda x:\mathsf{Nat}.\ \mathsf{succ}\ x)\ \mathsf{zero}
\\
&\qquad\longrightarrow
  \mathsf{succ}\ \mathsf{zero}.
\end{aligned}
\]

\section{Preservation Check}

The new evaluation rule preserves types.
Suppose
\[
\Gamma \vdash \mathbf{let}\ X = T\ \mathbf{in}\ t : [X \mapsto T]S.
\]
By inversion of \textsc{T-TLet},
\[
\Gamma \vdash T :: K
\qquad\text{and}\qquad
\Gamma, X :: K = T \vdash t : S.
\]
By the type-definition substitution lemma,
\[
\Gamma \vdash [X \mapsto T]t : [X \mapsto T]S.
\]
This is exactly the type of the reduct given by \textsc{E-TLet}.

Progress is also unchanged: a term of the form
$\mathbf{let}\ X = T\ \mathbf{in}\ t$ is not a value and can always take the
\textsc{E-TLet} step.

\end{document}
